\documentclass[12pt]{article}
\usepackage[utf8]{inputenc}
\usepackage[T1]{fontenc}
\usepackage{hyperref}
\title{Pinyin Language Resources for K-12 Curriculum Development}
\author{Compiled by Claude Sonnet 5.5 (claude-sonnet-5-5) for Andr\'as Kornai}
\date{}
\begin{document}
\maketitle

\section{Introduction}
Pinyin (ISO 639-3 pny; Glottocode piny1238) is a Grassfields (Atlantic-Congo) language of the Northwest Region of Cameroon, spoken by roughly 25,000 people (Mezam Division, near Bamenda). Other names include Bapinyi, Pelimpo and Menka. It is unrelated to Hanyu Pinyin, the romanization of Chinese. Searches found few resources.

\section{Primary Educational Resources}
No primers or textbooks were found in the searches. One source says the language is written in Latin script and used in schools, literature and media, but this was not corroborated.

\section{Linguistic and Reference Materials}
\begin{itemize}
\item Pinyin language (Wikipedia): \url{https://en.wikipedia.org/wiki/Pinyin_language}
\item Mbah and Anderson 2005, Pinyin orthography guide (CABTA; the Phoible source record): \url{https://phoible.org/sources/1632_mbah_pny}
\item Glottolog reference record for the orthography guide: \url{https://glottolog.org/resource/reference/id/50858}
\item IANA registration note on the language subtag: \url{https://www.iana.org/assignments/lang-subtags-templates/pny.txt}
\item Omniglot page (search result; mainly about Hanyu Pinyin, so possibly not relevant): \url{https://omniglot.com/writing/pinyin.htm}
\end{itemize}

\section{Religious and Community Resources}
\begin{itemize}
\item Joshua Project, Pinyin / Bapinyi in Cameroon: \url{https://joshuaproject.net/people_groups/14402}
\end{itemize}

\section{Digital and Audio Resources}
No audio, video or software resources were found.

\section{Curriculum Development Context}
A published orthography guide (2005) exists, which is a basis for materials. No dictionary or school programme was found. Related Grassfields languages (Bafut, Awing, Ngemba) have more developed literacy work that could serve as models.
\end{document}